\documentclass[12pt]{ctexart}   % 正文字体大小为12pt，中文支持

\usepackage{lipsum}             % 生成随机文字

% 页面大小与页边距设置
\usepackage{geometry}   
\geometry{a4paper, 
    left=2.54cm, 
    right=2.54cm, 
    top=3.18cm, 
    bottom=3.18cm
}

\linespread{1.5}        % 行距调整

% 页眉页脚设置
\usepackage{fancyhdr}
\pagestyle{fancy}
\fancyhf{}                      % 清空所有页眉和页脚区域的当前设置，避免继承默认的或其他样式残留的设置
\setlength{\headheight}{15pt}   % 设置页眉区域的高度为 15 pt
\addtolength{\topmargin}{-2pt}  % 调整页面上边距（\topmargin）的大小，将其减少 2 pt
\fancyfoot[C]{\thepage}         % 页脚居中显示页码
\fancyhead[L]{实验报告标题}       % 页眉左侧显示“实验报告标题”
\fancyhead[R]{作者1}             % 页眉右侧显示“作者1”


% 设置“封面”（包含标题、作者、日期等）
\title{\bfseries 实验报告标题}     % 标题
\author{作者1 \and 作者2}         % 作者
\date{\today}                    % 日期

\begin{document}
\maketitle      % 生成封面

\tableofcontents


\section{引言}

\lipsum[2-3]


\section{模型与实验设计}

\subsection{模型介绍}

\lipsum[4-5]


\subsection{实验设计}

\subsubsection{实验系统设计}

\lipsum[3-4]

\subsubsection{实验方法设计}

\lipsum[5-6]


\section{数据分析}

\subsection{关于XXX的分析}

\lipsum[1-3]

\subsection{关于YYY的分析}

\lipsum[7-8]


\section{结论与展望}

\lipsum[9]



\end{document}


